\documentclass[11pt,letterpaper]{article}
\usepackage[margin=1in]{geometry}
\usepackage[T1]{fontenc}
\usepackage{amsmath,amsthm,booktabs,array,longtable,graphicx}
\usepackage{newtxtext,newtxmath}
\usepackage[round,authoryear]{natbib}
\usepackage{setspace,xr-hyper,xurl}
\usepackage[hidelinks]{hyperref}
\usepackage{fancyhdr}
\onehalfspacing\setlength{\emergencystretch}{2em}
\newtheorem{theorem}{Theorem}[section]
\newtheorem{proposition}[theorem]{Proposition}
\newtheorem{lemma}[theorem]{Lemma}
\newtheorem{corollary}[theorem]{Corollary}
\theoremstyle{definition}\newtheorem{remark}[theorem]{Remark}
\newcommand{\R}{\mathbb R}\newcommand{\E}{\mathbb E}
\newcommand{\argmax}{\operatorname*{arg\,max}}
\newcommand{\argmin}{\operatorname*{arg\,min}}
\pagestyle{fancy}\fancyhf{}\fancyhead[L]{\small Finite-Catalog Resource Allocation}
\fancyhead[R]{\small Anonymous}\fancyfoot[C]{\thepage}\setlength{\headheight}{14pt}

% BEGIN retained/current source: revisions/or-r52-resource-path-20260925/generated/metrics.tex
\newcommand{\FineMaxGap}{0.000500}
\newcommand{\FineMinTime}{3.523}
\newcommand{\FineMaxTime}{19.762}
\newcommand{\FineMaxRelative}{1.63\%}
\newcommand{\FineMinPacked}{0.216}
\newcommand{\FineMaxPacked}{6.628}
\newcommand{\FineMaxDenBits}{553}
\newcommand{\AllResourceExact}{31}
\newcommand{\GroupNOneComplete}{2}
\newcommand{\MIPMaxSeconds}{0.038}
\newcommand{\MIPMaxGap}{$5.55\,10^{-17}$}

% END source


% BEGIN retained/current source: revisions/or-r54-exact-price-path-20260926/generated/metrics.tex
\newcommand{\PriceExact}{10}
\newcommand{\PriceTolerance}{19}
\newcommand{\PriceLimited}{1}
\newcommand{\PriceChecked}{30}
\newcommand{\PriceMinBytes}{541}
\newcommand{\PriceMaxBytes}{5895}
\newcommand{\AnytimeExact}{10}
\newcommand{\ForcedAnchors}{45}
\newcommand{\HardnessSubsets}{1692}
\newcommand{\HardnessBooks}{3004}

% END source

\hypersetup{pdftitle={Finite-Catalog Resource Allocation: A Tariff-Sensitive Complexity Frontier},pdfauthor={Anonymous}}
\externaldocument{.build/r64/main_refs}[main.pdf]
\externaldocument{.build/r64/electronic_companion_refs}[electronic_companion.pdf]
\begin{document}
\begin{center}{\Large\bfseries Response to the R63 Referee Report}\par\medskip
Finite-Catalog Resource Allocation: A Tariff-Sensitive Complexity Frontier\par
Revision R64 for \emph{Operations Research}\par September 27, 2026\end{center}

We thank the referee for the detailed independent examination of the mathematical results, source code, and frozen computational evidence. This revision addresses each of the five required revisions and the accompanying editorial comments. It retains the exact selected-boundary representation, original-model parameterized lower bound, tariff-exception algorithm, uniform-fee frontier, original-policy transfer theorem, conservation and approximation results, and same-book price characterization. No new computational campaign is substituted for the frozen study. The previous manuscript, proofs, response, execution sources, and failure records remain preserved.

\section*{1. Certificate-semantic closure}
\paragraph{Standard-fee selection (Required Revision 1(a)).}
The referee correctly distinguishes exact optimization conditional on a supplied standard from verification of the parameterization used to report the exception count. The independent tariff checker now recomputes frequencies from the original rational fee vector, selects the smallest most frequent fee, and rejects any different declared standard. It then retains every previous check of relative exception inventory, complete subset coverage, Bellman states, and original-policy recovery. This closes the claimed parameterization rather than narrowing the certificate to a conditional optimization statement. The main text preceding Theorem~\ref{thm:tariff60} explains that the recurrence is exact for any declared standard when its relative exceptions are fully enumerated; the mode minimizes its parameter, while the smallest-mode convention only fixes deterministic ties.

\paragraph{Optimal sparse projection and minimum allowance (Required Revision 1(b)).}
The independent projection checker reconstructs all retained windows at the claimed allowance. For each it computes the direct absolute-deviation sum at its lower median, rather than importing the optimizer's prefix-sum routine. It checks the minimizing window, deterministic ties, median, complete projected vector, relative exceptions, and minimum distortion. It also binds the inner exact tariff certificate to that projected vector. To certify minimality, it verifies that the selected allowance meets the total-distortion tolerance and that the preceding allowance fails it. Monotonicity then excludes every smaller allowance. The zero-allowance case is treated separately. The companion's Section~\ref{sec:semantic64} states the algorithm, its rational-operation count, and the precise scope of minimality. The compatibility verifier now delegates to this same independent verification path.

\paragraph{Adversarial tests.}
The new suite creates complete, internally consistent certificates carrying the disputed false metadata, rather than merely damaging a hash. It includes a nonoptimal retained window, an incorrect median, and a nonminimal selected allowance, as requested, as well as nonmodal and incorrect tied-mode standards. An isolated execution of the archived R63 checker accepts all ten adversarial cases; the revised checker rejects all ten. The suite also accepts 577 valid certificates and checks 288 projections against exhaustive retained-subset enumeration. A source-import check confirms that the production checker closure does not import the optimization modules. The machine-readable test records identify each adversary and the rejection reason.

The complete frozen inventory also replays successfully with the strengthened checkers: 405 requests remain inventoried, 342 two-sided rational certificates and 47 numerical-solver lower policies verify, and 16 requests remain without certificates. Sixteen certificate inputs retain their different but exactly equivalent rational serialization. All 45 later-successful replays retain their original failed or incomplete checking classification. The revision does not modify these empirical outcomes.

\section*{2. Exact tractability versus additive transfer}
The abstract now says that the exact sparse-exception solution is transferred to near-standard tariffs through an original-policy additive certificate. The introduction and conclusion use the same distinction. The exact fixed-parameter theorem remains unchanged: it concerns common linear rewards, free preliminary service, and the actual count of exceptions to a standard fee. For general near-standard original fees, the method exactly solves a tractable surrogate and evaluates its original-feasible policy under the original charges. The resulting interval is additive, not an unproved exact fixed-parameter algorithm for the original fees.

The full transfer theorem and its proof remain in Section~\ref{sec:robust63}. In particular, we retain $\bar m=\min\{m,N\}$ in the bound $2\bar m\|\delta\|_\infty$. The physical feasibility guarantees and the strict-margin book-stability statement are unchanged. We now state explicitly that the implementation does not compute a second-best-book margin: that condition is a theoretical diagnostic, not a reported stability certificate.

\section*{3. Antecedents and contribution hierarchy}
The related-work discussion now credits least trimmed absolute-deviation location and its consecutive-window/median construction, citing Tableman (1994) and Zioutas et al. (2015). It also identifies the fee comparison as deterministic objective-coefficient sensitivity analysis, with Xu and Burer (2017) as a broader optimal-value sensitivity reference. We do not claim the median, consecutive windows, or the elementary objective comparison as a new general clustering or sensitivity algorithm.

The model-specific contribution remains prominent: the exact tariff frontier identifies a tractable surrogate class, and its exact solution transfers to an original feasible policy with a certified original-fee value interval. The projection chooses the surrogate within a precise distortion criterion. This is neither statistical estimation of fees nor optimization over an unknown distribution. The projection proposition is labeled specialized sparse-tariff preprocessing; its proof remains available in full. These attributions clarify the supporting machinery without changing the scope or strength of the principal results.

\section*{4. Projection-quality evidence from the frozen records}
Table~\ref{tab:projection64} adds the requested compact evidence. Panel A reports catalog size, selected projection allowance, actual exceptions relative to the projected standard, the inner solver's modal exception count, both one-sided error budgets, the certified original-fee interval width, and exactly identified or bounded regret. Panel B reports the original optimizer and checker times and both uncompressed and gzip-compressed proof bytes. All requests have tolerance $1/500$ and the original six-second total allowance; these common quantities are stated in the caption.

The table includes all twelve adversarial robust-tariff requests, not only the seven that return certificates. Four are mathematically inapplicable and one encounters the engineering guard. Among the seven certified policies, five have exactly zero realized regret established by closed rational comparator intervals. The two diffuse near-standard cases have regret intervals $[0,8\times10^{-5}]$ and $[0,19\times10^{-5}]$, respectively; we do not describe those bounds as measured exact regret. Their zero-exception surrogates yield informative original-policy intervals even though their original tariffs have many exceptional commands.

The table-generation program reads frozen cases, raw records, and certificates and invokes only independent verification, never an optimizer. It records all source, row, and certificate hashes. Comparator proofs used after the study to identify regret retain their original recorded target and checking outcomes in the machine-readable output. No new timing observation enters the primary study, and no failed row is replaced.

\section*{5. Original timed checking and later integrity audits}
The primary computational captions now explicitly distinguish original target attainment from later external-input binding and semantic replay. An original success requires the originally assigned end-to-end deadline and original checker outcome. R63 binding and R64 semantic replay are later integrity audits; neither promotes an originally late or failed request. The independent root-support panel is separately identified as a post-study diagnostic, not another observation from the original timing environment. Its classification still requires support by one active original book, not a convex mixture of different books.

The complete 405-request inventory, 16 no-certificate requests, 45 later-replayed failures, engineering guards, and numerical-versus-rational evidence classes remain preserved. The table also retains probe and restart costs for the executed hybrid. We make no comparison with an ex post oracle selector. Numerical SCIP global bounds remain numerical; exact checking of a selected book certifies only the lower policy.

\section*{6. Remaining technical and editorial comments}
\paragraph{Notation, ties, tolerance, and guards (comments 2--5, 7, 12).}
The projection discussion distinguishes its allowed exception budget, the actual count relative to its projected standard, and the count relative to the inner solver's modal standard. These counts satisfy $d_{\rm tariff}\leq d_{\rm proj}\leq d_{\rm allow}$ and need not coincide. The implementation guard bounds the count used by the inner exact algorithm; it is not another mathematical parameter restriction. The tolerance controls total absolute fee distortion, a sufficient and possibly conservative condition for the final value width. Minimum allowance therefore does not mean minimum realized regret or minimum running time. Mode ties, conditional exactness, and the truncated command allowance are addressed explicitly in the article and companion.

\paragraph{Stability, projection, frontiers, and root support (comments 6, 8--10).}
The uncomputed second-best margin is identified as theoretical. The projection is supporting preprocessing, not another principal algorithmic novelty. The uniform-fee result continues to assert monotonicity of optimal cardinality, not nesting of installed sets. Root support continues to mean one active original book capable of implementing the promise.

\paragraph{Costs, evidence classes, and application scope (comments 11, 13--16, 18--19).}
The original hybrid costs, numerical upper bounds, adverse deficit results, and all missing phases are retained. The maximum proof sizes over the expanded panel are reported together as 5,572,345 uncompressed bytes and 214,897 compressed bytes; these are maxima, not a claim that checking is negligible. The projection table gives paired raw/compressed sizes for each available proof. The service-credit interpretation remains a stylized theory model, and the study remains constructed rather than field calibrated. Integrity receipts do not establish novelty, prove the theorems, or imply editorial acceptance.

\paragraph{Presentation, preservation, and ancestry (comments 1, 17, 20).}
The new manuscript, correspondence, and current repository instructions use \emph{Operations Research}. The article retains its main representation, lower bound, tariff theorem, robust-transfer theorem, and computational denominators. The additional checker details are placed in the companion. Previous readers and derivations are retained byte-for-byte in the history and archival snapshot, with a new preservation map. The revision descends from the reviewed author commit; the referee report is pinned by immutable commit and blob, not merged as a scientific parent. The final published checkout is separately verified and its exact commit is identified in the execution receipt.

We appreciate the referee's precise distinctions between mathematical guarantees, implementation claims, and empirical timing evidence. The revised manuscript preserves the tariff-sensitive complexity frontier while making its supporting certificates and empirical interpretation independently auditable.
\end{document}
